% =====================================================================
% experimento_binario_capacidad_costo_fixed.tex
% ---------------------------------------------------------------------
% Sección autocontenida: Experimento "binario_capacidad_20260531".
% Se integra en Capitulo_Experimentos.tex via \input{}.
% Sin preámbulo ni \begin{document}. Compilable con pdflatex >= 2 pases.
% =====================================================================

% Macro auxiliar de andamiaje: si ya existe en el documento padre se
% omite silenciosamente. Aquí se define de forma segura.
\providecommand{\TODO}[1]{\textcolor{rojomalo}{\textbf{[TODO: #1]}}}

% Macro de violación de capacidad (alias de notación).
\providecommand{\viol}{\mathrm{viol}}

% =====================================================================
\section{Experimento 9: binario\_capacidad\_20260531 (segundo approach con costo corregido)}
\label{sec:exp-binario-capacidad}
% =====================================================================

% =====================================================================
\subsection{Hipótesis y motivación}
% =====================================================================

El Experimento~8 (\texttt{solo\_p\_inter\_20260531}, Sección~\ref{sec:exp-solo-p-inter})
estableció la primera línea de base bajo el evaluador greedy integrado de
forma nativa. Su selector es \emph{probabilístico}: en cada iteración,
el parámetro $p_{\text{inter}}$ determina con qué frecuencia se propone
el grupo INTER-ruta cuando la solución es factible; el parámetro
$\alpha_{\text{inter}} = 0.80$ hace lo propio en estado infactible.

El presente experimento sustituye el selector probabilístico por un
\textbf{selector binario determinista} basado en la violación de
capacidad: si la solución viola la restricción, se propone siempre el
grupo INTER-ruta (reparación); si es factible, siempre el grupo
INTRA-ruta (refinamiento). La consecuencia metodológica directa es que
$p_{\text{inter}}$ deja de ser un parámetro: el selector es determinista
y la única diferencia frente al Experimento~8 es la regla de selección.

La hipótesis central del experimento es la siguiente:

\begin{quote}
\textit{Un selector completamente determinista que impone una
separación estricta entre las fases de reparación y refinamiento
puede igualar o superar al selector probabilístico calibrado
(\texttt{solo\_p\_inter}), porque elimina el ruido de la decisión
aleatoria y concentra el presupuesto de movimientos INTER en los
momentos en que son necesarios y el presupuesto INTRA cuando la
solución es factible, sin desperdiciar iteraciones en combinaciones
subóptimas de grupo/estado.}
\end{quote}

Si la comparativa entre el Experimento~8 y el Experimento~9 muestra
gaps similares, el resultado indica que la señal binaria de
factibilidad es suficiente para guiar la selección de grupo y que
$p_{\text{inter}}$ no aporta valor adicional una vez que el evaluador
es correcto. Si el selector binario resulta inferior, el siguiente paso
natural es explorar selectores adaptativos o híbridos (por ejemplo,
binario con suavizado durante la fase factible).

% =====================================================================
\subsection{Diseño del approach}
% =====================================================================

El approach \texttt{binario\_capacidad} comparte con el Experimento~8
todos los componentes excepto el selector de grupo de operadores:

\begin{itemize}
\item Mismo evaluador de costo greedy nativo
  (\texttt{metacarp/evaluador\_costo.py}).
\item Mismo conjunto de 9 operadores
  (\texttt{OPERADORES\_POPULARES}: 3 intra + 6 inter).
\item Mismo $\lambda$ de penalización de capacidad por defecto
  (\texttt{lambda\_capacidad=None}, $\approx 10\times$ mediana del
  costo de arco, con mínimo absoluto de 10).
\item Mismos parámetros instance-aware (\texttt{None}) en todas las MH.
\item Sin Path Relinking, sin kick reactivo, sin AOS, sin presupuesto
  extendido.
\item Sin semilla fija: la variabilidad entre repeticiones es
  puramente estocástica.
\end{itemize}

\noindent Lo que cambia es el selector:

\begin{itemize}
\item \textbf{Selector binario determinista}
  (\texttt{seleccionar\_grupo\_strict}, módulo
  \texttt{metacarp/strict\_intra\_inter\_20260524.py}): decide el
  grupo de operadores según el signo de la violación de capacidad
  $\viol$ de la solución actual.
  \begin{itemize}
  \item Si $\viol > 10^{-12}$: proponer siempre el grupo
    \textbf{INTER-ruta} (reparación).
  \item Si $\viol \le 10^{-12}$: proponer siempre el grupo
    \textbf{INTRA-ruta} (refinamiento).
  \end{itemize}
\item El selector \textbf{no consume ningún} \texttt{rng.random()} por
  iteración: es estrictamente determinista una vez fijada la
  violación. Ignora \texttt{p\_inter} y \texttt{alpha\_inter},
  absorbiéndolos en \texttt{**\_ignorado}.
\item Dentro del grupo elegido, \texttt{generar\_vecino} selecciona el
  operador de forma \textbf{uniforme} (\texttt{pesos\_operadores=None}),
  igual que en el Experimento~8.
\end{itemize}

\medskip
\noindent\textbf{Kick reactivo desactivado.} El módulo
\texttt{strict\_intra\_inter\_20260524} expone también funciones de
kick inter-ruta (\texttt{aplicar\_kick\_labels},
\texttt{aplicar\_kick\_ids}) que el experimento original
\texttt{strict\_intra\_inter\_20260524} sí activaba. En el presente
approach el kick queda explícitamente desactivado
(\texttt{max\_iter\_sin\_mejora\_kick=None}) para aislar el efecto del
selector binario puro, sin la capa de diversificación. Esta decisión
mantiene la comparabilidad directa con el Experimento~8.

\medskip
\noindent\textbf{Instalación del selector vía monkey-patch.}
El selector binario se instala en el módulo de cada MH mediante la
función \texttt{aplicar\_patch\_selector} del script
\texttt{scripts/\_strict\_intra\_inter\_20260524\_common.py}. La
función reasigna el atributo
\texttt{seleccionar\_grupo\_operadores\_inter\_intra} en el namespace
del módulo de la MH por la versión binaria estricta, de modo que la
MH la invoque sin modificación alguna en su lógica interna. El patch
se aplica al inicio de cada worker (proceso hijo del
\texttt{ProcessPoolExecutor}), porque cada proceso hijo importa el
módulo de la MH desde cero.

La diferencia metodológica respecto al Experimento~8 se resume en la
Tabla~\ref{tab:comparativa-approaches-binario-capacidad}.

\begin{table}[h]
\centering
\caption{Diferencias entre el Experimento~8 (\texttt{solo\_p\_inter})
         y el Experimento~9 (\texttt{binario\_capacidad}).}
\label{tab:comparativa-approaches-binario-capacidad}
\small
\begin{tabular}{lll}
\toprule
Componente & Exp.~8 \texttt{solo\_p\_inter} & Exp.~9 \texttt{binario\_capacidad} \\
\midrule
Selector de grupo        & Probabilístico ($p_{\text{inter}}$)   & \textbf{Binario determinista} \\
Parámetro $p_{\text{inter}}$ & Calibrado en $\{0.1,\ldots,0.5\}$ & \textbf{No existe} \\
Consume \texttt{rng} por iter. & 1 llamada                        & \textbf{0 llamadas} \\
Monkey-patch selector    & No                                    & \textbf{Sí (\texttt{aplicar\_patch\_selector})} \\
Kick reactivo            & No                                    & No (aislamiento) \\
Configs por MH en grid   & 15 ($5 \times 3$)                     & \textbf{3 (solo libre)} \\
Corridas totales (Fase 1)& 5\,175                                & \textbf{1\,035} \\
Conjunto operadores      & 9 (\texttt{POPULARES})                & 9 (\texttt{POPULARES}) \\
$\lambda$                & Default (None)                        & Default (None) \\
Parámetros MH            & Instance-aware (None)                 & Instance-aware (None) \\
Evaluador de costo       & Greedy nativo                         & Greedy nativo \\
\bottomrule
\end{tabular}
\end{table}

% =====================================================================
\subsection{Selector binario: pseudocódigo}
\label{subsec:selector-binario-capacidad}
% =====================================================================

El Algoritmo~\ref{alg:selector-binario-capacidad} formaliza el
mecanismo. A diferencia del selector probabilístico
(Algoritmo~\ref{alg:selector-solo-p-inter}), no hay ninguna extracción
de número aleatorio: la rama elegida depende exclusivamente del signo
de $\viol$.

\begin{algorithm}[h]
\caption{Selector binario determinista basado en capacidad
         (\texttt{seleccionar\_grupo\_strict}).}
\label{alg:selector-binario-capacidad}
\begin{algorithmic}[1]
\Require $\viol \in \mathbb{R}_{\ge 0}$ violación de capacidad actual;
  conjuntos $\mathcal{O}_{\text{intra}},\,\mathcal{O}_{\text{inter}},\,\mathcal{O}_{\text{fb}}$
  (fallback).
  Los argumentos \texttt{rng}, \texttt{p\_inter}, \texttt{alpha\_inter}
  se reciben y se ignoran (\texttt{**\_ignorado}).
\Ensure Grupo de operadores $\mathcal{O}$ a usar en esta iteración.
\If{$\viol > 10^{-12}$ \textbf{y} $\mathcal{O}_{\text{inter}} \neq \emptyset$}
  \State \Return $\mathcal{O}_{\text{inter}}$, $\mathit{hubo\_violacion}{=}\mathrm{True}$
  \Comment{Estado infactible: reparación forzada}
\ElsIf{$\mathcal{O}_{\text{intra}} \neq \emptyset$}
  \State \Return $\mathcal{O}_{\text{intra}}$, $\mathit{hubo\_violacion}{=}(\viol > 10^{-12})$
  \Comment{Estado factible: refinamiento}
\ElsIf{$\mathcal{O}_{\text{inter}} \neq \emptyset$}
  \State \Return $\mathcal{O}_{\text{inter}}$, $\mathit{hubo\_violacion}{=}(\viol > 10^{-12})$
  \Comment{Intra vacío: caída al inter}
\Else
  \State \Return $\mathcal{O}_{\text{fb}}$, $\mathit{hubo\_violacion}{=}(\viol > 10^{-12})$
  \Comment{Caso degenerado: fallback completo}
\EndIf
\end{algorithmic}
\end{algorithm}

El valor de retorno es la tupla
$(\text{grupo\_operadores},\,\mathit{hubo\_violacion})$, con la misma
semántica que el helper canónico
\texttt{seleccionar\_grupo\_operadores\_inter\_intra}: el flag
$\mathit{hubo\_violacion}$ lo consumen TS simple y RTS para sus
contadores internos de iteraciones con violación de capacidad.

Un aspecto relevante para la reproducibilidad: dado que el selector no
consume ningún \texttt{rng.random()}, las secuencias pseudoaleatorias
de dos corridas idénticas (misma instancia, misma semilla, misma MH)
diferirán de las del Experimento~8 en exactamente \textbf{una llamada
por iteración} (la que en el Experimento~8 realizaba el selector
probabilístico). Esto hace que los resultados de ambos approaches no
sean directamente comparables corrida a corrida, pero sí agregadamente
(gap medio sobre múltiples repeticiones sin semilla fija).

% =====================================================================
\subsection{Configuración fija canónica}
\label{subsec:config-fija-binario-capacidad}
% =====================================================================

Este approach NO calibra: usa la \textbf{configuración canónica fija} por
MH ---la misma que el Experimento~8--- y cambia \emph{únicamente} el
selector. El selector binario es determinista y \emph{no tiene el
parámetro $p_{\text{inter}}$}, así que la única diferencia operativa frente
al Experimento~8 es la regla de selección de grupo. Se ejecuta una sola
configuración por MH con \textbf{5 repeticiones} $\times$ 23 instancias
(sin semilla; la variabilidad es estocástica), y los CSVs completos (con
todas las columnas de operadores y deadheading) se generan en
\texttt{final/}. La config usada se registra en \texttt{config\_fija.json}.

\begin{table}[h]
\centering
\caption{Configuración canónica fija por metaheurística en el approach
         \texttt{binario\_capacidad} (idéntica a la del Experimento~8; el
         selector binario ignora $p_{\text{inter}}$). El resto de
         parámetros se dejan instance-aware (\texttt{None}).}
\label{tab:grid-binario-capacidad}
\small
\begin{tabular}{llll}
\toprule
MH & Parámetro libre & Valor fijo & Resto notable \\
\midrule
SA           & \texttt{alpha} (enfriamiento)   & $0.90$  &
  $T_{\text{init}}, T_{\min}$: instance-aware; reheat acotado\textsuperscript{†} \\
TS simple    & \texttt{tabu\_tenure}            & $25$    &
  tamaño vecindario: 25 \\
RTS          & \texttt{factor\_aumento}         & $1.2$   &
  \texttt{factor\_reduccion}: $0.9$ \\
ABC simple   & \texttt{num\_fuentes}            & $30$    &
  límite abandono: instance-aware \\
Cuckoo       & \texttt{pa\_abandono}            & $0.15$  &
  \texttt{beta\_levy}: $1.5$; nidos: instance-aware \\
\bottomrule
\end{tabular}
\\[4pt]
\footnotesize
\textsuperscript{†}Reheat acotado SA: \texttt{patience=10},
\texttt{reheat\_factor=0.5}, \texttt{max\_reheats\_sin\_mejora=5}.
Véase la Nota~\ref{nota:reheat-acotado-binario-capacidad}.
\end{table}

\paragraph{Corridas totales.}
\begin{itemize}
\item Fase 1: $3\,\text{configs} \times 23\,\text{inst} \times
  3\,\text{reps} = 207$ corridas por MH
  $= 1{,}035$ corridas totales.
\item Fase 2: $23\,\text{inst} \times 5\,\text{reps} = 115$ corridas
  por MH $= 575$ corridas totales.
\item \textbf{Total}: 1{,}610 corridas entre ambas fases
  (frente a 5{,}750 del Experimento~8, una reducción de $\approx 72$\,\%).
\end{itemize}

% =====================================================================
\subsection{Nota técnica: reheat acotado en SA}
\label{nota:reheat-acotado-binario-capacidad}
% =====================================================================

El Recocido Simulado utiliza un mecanismo de reheat para salir de
mínimos locales cuando la temperatura cae demasiado rápido. Con el
reheat por defecto del proyecto (\texttt{max\_reheats\_sin\_mejora=0},
es decir, sin tope) y un valor de enfriamiento lento como
$\alpha = 0.99$, la temperatura se mantiene alta durante muchas
iteraciones, lo que puede alargar la corrida de forma impráctica.

Para garantizar que todas las corridas de SA terminen en un tiempo
razonable independientemente del valor de $\alpha$, el approach
\texttt{binario\_capacidad} fija el reheat \textbf{acotado}:

\begin{itemize}
\item \texttt{patience=10}: número de iteraciones sin mejora antes de
  activar el reheat.
\item \texttt{reheat\_factor=0.5}: la temperatura de reheat es la mitad
  de la temperatura actual.
\item \texttt{max\_reheats\_sin\_mejora=5}: máximo de reheats
  consecutivos sin mejora del mejor global antes de detener la búsqueda.
\end{itemize}

Esta configuración es la misma que la validada en el programa
experimental del primer ciclo y garantiza que cada corrida de SA
converja en un tiempo comparable al de las demás metaheurísticas,
con independencia del valor de $\alpha$ explorado en el grid.

% =====================================================================
\subsection{Configuración y parámetros instance-aware}
% =====================================================================

La Tabla~\ref{tab:config-binario-capacidad} resume la configuración
completa del approach y la compara con el Experimento~8 y con el
experimento original \texttt{strict\_intra\_inter\_20260524} (que sí
activaba el kick reactivo y usaba el subconjunto reducido de 5
operadores).

\begin{table}[h]
\centering
\caption{Comparativa de configuración entre el experimento original
         \texttt{strict\_intra\_inter\_20260524}, el Experimento~8
         (\texttt{solo\_p\_inter}) y el Experimento~9
         (\texttt{binario\_capacidad}).}
\label{tab:config-binario-capacidad}
\small
\begin{tabular}{llll}
\toprule
Componente &
  \texttt{strict\_intra\_inter} &
  Exp.~8 \texttt{solo\_p\_inter} &
  Exp.~9 \texttt{binario\_capacidad} \\
\midrule
Selector de grupo    & Binario strict       & Probabilístico ($p_{\text{inter}}$) & \textbf{Binario strict} \\
Kick reactivo        & Sí                   & No                                  & \textbf{No} \\
Conjunto operadores  & 5 (\texttt{STRICT\_5}) & 9 (\texttt{POPULARES})            & \textbf{9 (\texttt{POPULARES})} \\
Evaluador de costo   & Greedy (patch)       & Greedy nativo                       & Greedy nativo \\
Monkey-patch         & Selector + kick      & No                                  & \textbf{Solo selector} \\
$p_{\text{inter}}$   & Ignorado             & Calibrado                           & \textbf{Ignorado} \\
$\lambda$            & Default (None)       & Default (None)                      & Default (None) \\
Reheat SA            & Default (sin tope)   & Acotado (\texttt{max\_reheats=5})   & Acotado (\texttt{max\_reheats=5}) \\
Parada otras MH      & Instance-aware       & Instance-aware                      & Instance-aware \\
Grid configs por MH  & ---                  & 15 ($5 \times 3$)                   & \textbf{3} \\
\bottomrule
\end{tabular}
\end{table}

\medskip
\noindent\textbf{Parámetros instance-aware.} Idénticos al Experimento~8:
$T_{\text{init}}$ y $T_{\min}$ en SA; número de iteraciones y tenure
bounds en TS/RTS; límite de abandono en ABC; número de nidos y pasos
de Lévy en Cuckoo. Todos se pasan como \texttt{None} al wrapper
\texttt{*\_desde\_instancia} de cada MH.

% =====================================================================
\subsection{Reproducibilidad: scripts e infraestructura}
% =====================================================================

El experimento se organiza en tres niveles de scripts, en paralelo con
la estructura del Experimento~8:

\begin{itemize}
\item \textbf{Orquestador común}:
  \texttt{scripts/\_binario\_capacidad\_20260531\_common.py} ---
  corre la configuración canónica fija (sin grid), la construcción de
  kwargs por MH (incluyendo la instalación del selector binario vía
  \texttt{aplicar\_patch\_selector}), la consolidación de CSV parciales
  y la CLI (\texttt{--smoke}, \texttt{--reps}, \texttt{--workers}).
\item \textbf{Runners por MH}:
  \texttt{scripts/run\_sa\_binario\_capacidad\_20260531.py},
  \texttt{run\_tabu\_simple\_binario\_capacidad\_20260531.py},
  \texttt{run\_tabu\_reactiva\_binario\_capacidad\_20260531.py},
  \texttt{run\_abc\_simple\_binario\_capacidad\_20260531.py},
  \texttt{run\_cuckoo\_binario\_capacidad\_20260531.py} ---
  cada uno fija su MH y llama al orquestador.
\item \textbf{Lanzador master}:
  \texttt{scripts/run\_all\_binario\_capacidad\_20260531.sh} ---
  lanza las cinco MH en secuencia (una por vez, saturando la CPU
  internamente con \texttt{ProcessPoolExecutor}), con redirección de
  logs a \texttt{logs/}.
\end{itemize}

La salida se escribe en:
\begin{center}
\texttt{experimentos\_costo\_fixed/<mh>\_binario\_capacidad\_<YYYYMMDD-HHMM>/}
\end{center}

Con la siguiente estructura de carpetas:

\begin{itemize}
\item \texttt{config\_fija.json} --- la configuración canónica usada
  (MH, knob libre y su valor), para trazabilidad.
\item \texttt{final/} --- CSVs de la corrida (5 repeticiones, columnas
  completas), uno por instancia con el patrón
  \texttt{<mh>\_binario\_capacidad\_<instancia>.csv}.
  El subdirectorio \texttt{final/\_partials/} contiene los parciales
  antes de consolidar.
\end{itemize}

\medskip
\noindent\textbf{Columnas del CSV.} Idénticas a las del Experimento~8
y al núcleo del proyecto: \texttt{gap\_bks\_porcentaje},
\texttt{mejor\_costo}, \texttt{costo\_total\_desde\_reporte},
\texttt{reporte\_detalle\_deadheading}, y los 36 contadores de
operadores ($9\text{ operadores} \times 4\text{ categorías}$:
propuesto, aceptado, mejoró, trayectoria\_mejor). El campo
\texttt{extra\_csv} registra
\texttt{selector="binario\_estricto\_capacidad"} y
\texttt{approach="binario\_capacidad"} para trazabilidad.

% =====================================================================
\subsection{Resultados}
% =====================================================================

\begin{observacion}[Corrida completada]
El experimento \texttt{binario\_capacidad\_20260531} se ejecutó el
31 de mayo de 2026 sobre las 23 instancias, con grid search en dos
fases (3 repeticiones de calibración y 5 de confirmación). Los valores
reportados corresponden a la fase~2 con la configuración ganadora de
cada MH.
\end{observacion}

\paragraph{Fuentes de datos.}
Cada número de la Tabla~\ref{tab:binario-capacidad-resultados} se
extrae de los siguientes archivos generados por el orquestador:

\begin{itemize}
\item \textbf{Parámetro libre* y valor*}: la config canónica fija de
  \texttt{config\_fija.json}.
\item \textbf{Gap medio}: promedio de la columna
  \texttt{gap\_bks\_porcentaje} en los CSV de
  \texttt{final/<mh>\_binario\_capacidad\_<instancia>.csv} sobre las
  23 instancias $\times$ 5 repeticiones.
\item \textbf{$N$ instancias en BKS}: número de instancias para las
  que \texttt{gap\_bks\_porcentaje} $< 0.01$\,\% en al menos una de
  las 5 repeticiones.
\end{itemize}

\begin{table}[h]
\centering
\caption{Resultados del experimento \texttt{binario\_capacidad\_20260531}
         (config fija, 5 reps $\times$ 23 instancias por MH). El gap se mide
         respecto al BKS con el evaluador greedy nativo.}
\label{tab:binario-capacidad-resultados}
\small
\begin{tabular}{lccccc}
\toprule
MH & Parámetro libre* & Valor* &
  Gap medio (\%) & Mediana (\%) & $N$ inst.\ en BKS \\
\midrule
SA           & \texttt{alpha}           & 0.90 &
  21.04 & 19.93 & 0/23 \\
TS simple    & \texttt{tabu\_tenure}    & 25   &
  19.14 & 18.00 & 1/23 \\
RTS          & \texttt{factor\_aumento} & 1.2  &
  14.82 & 13.13 & 1/23 \\
ABC simple   & \texttt{num\_fuentes}    & 30   &
  15.39 & 15.00 & 2/23 \\
Cuckoo       & \texttt{pa\_abandono}    & 0.15 &
  21.30 & 20.75 & 0/23 \\
\midrule
\textbf{Media} & --- & --- & \textbf{18.34} & --- & 4/115 \\
\bottomrule
\end{tabular}
\\[4pt]
\footnotesize
* Configuración canónica fija (el mismo valor por MH en los tres approaches;
leída de \texttt{config\_fija.json}). La fila \textbf{Media} agrega los aciertos
de BKS como pares (MH, instancia) sobre el total $5 \times 23$ por MH.
\end{table}

\paragraph{Comparativa Experimento 8 vs.\ Experimento 9.}
La Tabla~\ref{tab:comparativa-exp8-exp9} sitúa los resultados de ambos
approaches frente a frente por MH. El objetivo es cuantificar el efecto
de sustituir el selector probabilístico calibrado ($p_{\text{inter}}$)
por el selector binario determinista, con todo lo demás idéntico.

\begin{table}[h]
\centering
\caption{Gap medio (\%) por MH: Experimento~8 (\texttt{solo\_p\_inter},
         selector probabilístico con $p_{\text{inter}}$ calibrado)
         frente a Experimento~9 (\texttt{binario\_capacidad}, selector
         binario determinista). La columna $\Delta$ indica la diferencia
         Exp.~9 $-$ Exp.~8 en puntos porcentuales (negativo = Exp.~9
         mejor).}
\label{tab:comparativa-exp8-exp9}
\small
\begin{tabular}{lrrrl}
\toprule
MH &
  Exp.~8 \texttt{solo\_p\_inter} &
  Exp.~9 \texttt{binario\_capacidad} &
  $\Delta$ (pp) &
  Interpretación \\
\midrule
SA        & 3.21  & 21.04 & $+17.83$ & Exp.~8 mucho mejor \\
TS simple & 8.40  & 19.14 & $+10.74$ & Exp.~8 mucho mejor \\
RTS       & 4.89  & 14.82 & $+9.93$  & Exp.~8 mucho mejor \\
ABC       & 14.66 & 15.39 & $+0.73$  & equivalentes \\
Cuckoo    & 19.91 & 21.30 & $+1.39$  & equivalentes \\
\midrule
\textbf{Media} &
  \textbf{10.21} & \textbf{18.34} & $+8.13$ & Exp.~8 mejor \\
\bottomrule
\end{tabular}
\\[4pt]
\footnotesize
Los gaps del Exp.~8 provienen de
\texttt{final/<mh>\_solo\_p\_inter\_<instancia>.csv};
los del Exp.~9 de
\texttt{final/<mh>\_binario\_capacidad\_<instancia>.csv}.
\end{table}

\paragraph{Interpretación de los resultados.}
La comparativa es contundente: $\Delta > 0$ en las cinco MH (Exp.~8
mejor en todas) con una diferencia media de $+8.13$\,pp. El selector
binario determinista es \textbf{claramente inferior} al probabilístico
calibrado, y la brecha es dramática en las metaheurísticas de
búsqueda local: SA pasa de 3.21\,\% a 21.04\,\% ($+17.83$\,pp), TS de
8.40\,\% a 19.14\,\% y RTS de 4.89\,\% a 14.82\,\%. En las poblacionales
(ABC, Cuckoo) la diferencia es marginal ($+0.73$ y $+1.39$\,pp), pero
porque ya partían de gaps altos.

La causa es estructural: el selector binario \textbf{nunca propone
movimientos inter-ruta cuando la solución es factible}. Una vez que la
búsqueda alcanza una región factible —lo habitual en estas instancias
pequeñas— queda confinada a operadores intra-ruta y se estanca en un
óptimo local intra, incapaz de redistribuir tareas entre rutas. El
selector probabilístico del Experimento~8, con $p_{\text{inter}}$ alto
(0.4--0.5), mantiene un flujo constante de exploración inter-ruta
incluso en estado factible, que es justamente lo que permite cerrar el
gap. Esto refuerza el hallazgo del Experimento~8: bajo el costo
corregido, la exploración inter-ruta sostenida es el factor dominante.

En consecuencia, el selector binario \emph{puro} queda descartado como
base competitiva. El selector probabilístico $p_{\text{inter}}$ es la
base de referencia sobre la que conviene añadir los mecanismos de
intensificación de los siguientes approaches (Path Relinking en el
Experimento~10), sin perjuicio de que el selector binario pueda
recuperar competitividad combinado con un kick reactivo que reintroduzca
movimientos inter-ruta de forma periódica.

% =====================================================================
\subsection{Conclusión provisional del experimento 9}
% =====================================================================

Este experimento es el segundo approach del nuevo ciclo experimental
bajo el evaluador greedy integrado de forma nativa. Más allá de los gaps
obtenidos, sus contribuciones metodológicas son:

\begin{enumerate}
\item \textbf{Aislamiento del selector como variable experimental.}
  Al usar la MISMA configuración canónica por MH que el Experimento~8
  (mismo \texttt{alpha}/\texttt{tabu\_tenure}/\dots, operadores, $\lambda$,
  parámetros instance-aware, sin PR/kick/AOS) y cambiar \emph{solo} el
  selector, la comparativa Exp.~8 vs.\ Exp.~9 atribuye directamente la
  diferencia de gaps al tipo de selector: probabilístico
  ($p_{\text{inter}}$) frente a binario determinista.
\item \textbf{Separación estricta de las fases de búsqueda.}
  El selector binario implementa de forma explícita la intuición de
  que la reparación de la factibilidad y el refinamiento de la solución
  son tareas distintas que se benefician de operadores distintos.
  Cuantificar si esta intuición se traduce en gaps menores bajo el
  evaluador correcto es una contribución empírica al diseño de
  metaheurísticas para CARP.
\end{enumerate}

Los resultados (Tablas~\ref{tab:binario-capacidad-resultados}
y~\ref{tab:comparativa-exp8-exp9}) confirman que el selector binario
\emph{puro} no es competitivo bajo el costo corregido; su valor reaparece
solo combinado con intensificación (véase el Experimento~\ref{sec:exp-pr-aislado}).
